% =====================================================================
% CM Séance 1 — ADTs et Interfaces Java
% Structures de Données et Algorithmes Avancés — M1 — ESP/UCAD
% Dr. El Hadji Bassirou TOURÉ — 2025–2026
% Compilation : pdflatex (2 passes)
% =====================================================================
\documentclass[11pt,a4paper]{article}

\usepackage[utf8]{inputenc}
\usepackage[T1]{fontenc}
\usepackage{lmodern}
\usepackage[french]{babel}

\usepackage[margin=2.5cm]{geometry}
\usepackage{amsmath,amssymb,amsthm}
\usepackage{graphicx}
\graphicspath{{figures/}}
\usepackage{xcolor}
\usepackage{booktabs}
\usepackage{array}
\usepackage{enumitem}
\usepackage{fancyhdr}
\usepackage{longtable}
\usepackage{multirow}
\usepackage{pifont}
\usepackage{listings}
\usepackage{tcolorbox}
\tcbuselibrary{theorems,breakable,skins}
\usepackage{hyperref}

\hypersetup{colorlinks=true,linkcolor=blue!60!black,urlcolor=blue!60!black,
            citecolor=blue!60!black}

\newcommand{\cmark}{\ding{51}}
\newcommand{\xmark}{\ding{55}}
\newcommand{\bigO}{\mathcal{O}}
\newcommand{\bigOmega}{\Omega}
\newcommand{\bigTheta}{\Theta}

% ===== Boîtes =====
\newtcolorbox{objectifbox}[1][]{colback=blue!5,colframe=blue!50,
    fonttitle=\bfseries,title=#1,breakable}
\newtcolorbox{infobox}[1][]{colback=green!5,colframe=green!50!black,
    fonttitle=\bfseries,title=#1,breakable}
\newtcolorbox{attentionbox}[1][]{colback=orange!5,colframe=orange!70,
    fonttitle=\bfseries,title=#1,breakable}
\newtcolorbox{retenirbox}[1][]{colback=yellow!10,colframe=yellow!50!black,
    fonttitle=\bfseries,title=#1,breakable}
\newtcolorbox{prereqbox}[1][]{colback=cyan!5,colframe=cyan!50!black,
    fonttitle=\bfseries,title=#1,breakable}
\newtcolorbox{theorembox}[1][]{colback=purple!5,colframe=purple!50,
    fonttitle=\bfseries,title=#1,breakable}
\newtcolorbox{appliboxQ}[1][]{colback=gray!8,colframe=gray!60,
    fonttitle=\bfseries,title=#1,breakable}
\newtcolorbox{appliboxR}[1][]{colback=green!3,colframe=green!40!black,
    fonttitle=\bfseries\itshape,title=#1,breakable}
\newtcolorbox{exercicebox}[1][]{colback=white,colframe=gray!70,
    fonttitle=\bfseries,title=#1,breakable}
\newtcolorbox{lienbox}[1][]{colback=blue!3,colframe=blue!70!black,
    fonttitle=\bfseries\itshape,title=#1,breakable,
    before upper={\itshape}}

% ===== Théorèmes =====
\theoremstyle{definition}
\newtheorem{definition}{Définition}[section]
\newtheorem{theorem}{Théorème}[section]
\newtheorem{proposition}[theorem]{Proposition}
\theoremstyle{remark}
\newtheorem{example}{Exemple}[section]
\newtheorem{remark}{Remarque}[section]

% ===== Code =====
\lstset{
  language=Java,
  basicstyle=\ttfamily\small,
  backgroundcolor=\color{gray!8},
  frame=single,
  framerule=0.5pt,
  rulecolor=\color{gray!40},
  keywordstyle=\color{blue!70!black}\bfseries,
  commentstyle=\color{green!50!black}\itshape,
  stringstyle=\color{red!60!black},
  numbers=left,
  numberstyle=\tiny\color{gray},
  breaklines=true,
  breakatwhitespace=true,
  showstringspaces=false,
  tabsize=4,
  extendedchars=true,
  xleftmargin=4pt,
  xrightmargin=4pt,
  aboveskip=8pt,
  belowskip=8pt,
  literate={é}{{\'e}}1 {è}{{\`e}}1 {ê}{{\^e}}1 {à}{{\`a}}1 {ù}{{\`u}}1
           {ô}{{\^o}}1 {î}{{\^i}}1 {ç}{{\c{c}}}1 {û}{{\^u}}1 {â}{{\^a}}1
           {ë}{{\"e}}1 {ï}{{\"i}}1 {É}{{\'E}}1 {È}{{\`E}}1 {À}{{\`A}}1
           {Ç}{{\c{C}}}1 {Ê}{{\^E}}1 {œ}{{\oe}}1 {°}{{$^\circ$}}1
           {«}{{\guillemotleft}}1 {»}{{\guillemotright}}1
}

% ===== En-têtes =====
\pagestyle{fancy}
\fancyhf{}
\fancyhead[L]{\small S01 --- ADTs et Interfaces Java}
\fancyhead[R]{\small Structures de Données Avancées --- M1}
\fancyfoot[L]{\small Dr. El Hadji Bassirou TOURÉ}
\fancyfoot[C]{\thepage}
\fancyfoot[R]{\small 2025--2026}
\renewcommand{\headrulewidth}{0.4pt}
\renewcommand{\footrulewidth}{0.4pt}

\title{%
\vspace{-1cm}
\textbf{\LARGE Structures de Données et}\\
\textbf{\LARGE Algorithmes Avancés}\\[0.8cm]
\textbf{\huge Séance 1 --- ADTs et Interfaces Java}\\[0.3cm]
\textbf{\Large Comment séparer le contrat de l'implémentation ?}\\[0.5cm]
\large Master 1 --- 2025--2026}
\author{Dr. El Hadji Bassirou TOURÉ\\
Département Génie Informatique\\
Université Cheikh Anta Diop de Dakar}
\date{}

\begin{document}
\maketitle
\thispagestyle{empty}

\begin{prereqbox}[Navigation conceptuelle]
\textbf{Ce que vous savez déjà :} Java orienté objet (classes, héritage,
interfaces), et l'environnement de travail installé au TP-1 (Docker, Gradle,
JUnit).\\[4pt]
\textbf{La limite qu'on dépasse aujourd'hui :} vous savez écrire une classe
qui stocke des données. Mais dès qu'un programme grossit, le code client se
retrouve soudé à une décision technique --- « on utilise un tableau » ---
qu'on ne peut plus changer sans tout réécrire.\\[4pt]
\textbf{Ce que vous saurez faire à la fin :}
\begin{enumerate}[nosep]
    \item Énoncer un type abstrait de données (ADT) sous forme d'interface
          Java générique, indépendamment de toute implémentation.
    \item Implémenter List, Stack, Queue et Deque par tableau et par nœuds
          chaînés, et énoncer les compromis de chaque choix.
    \item Utiliser une pile pour vérifier un parenthésage et évaluer une
          expression postfixe.
    \item Choisir, pour un problème donné, l'ADT adapté avant de choisir la
          structure concrète.
\end{enumerate}
\textbf{Ouvre la voie vers :} S2, où l'on se donnera les outils mathématiques
pour \emph{prouver} qu'une implémentation est meilleure qu'une autre.
\end{prereqbox}

\tableofcontents
\newpage

% =====================================================================
\section{D'où l'on part}
% =====================================================================

En licence, vous avez appris à écrire des classes : un tableau ici, une liste
chaînée là, quelques méthodes pour ajouter et retirer des éléments. Cela
suffit tant que le programme tient dans un fichier. Dès qu'il grandit, une
difficulté apparaît : le code qui \emph{utilise} la structure connaît la
manière dont elle est construite, et il en dépend. Remplacer le tableau par
une liste chaînée oblige alors à corriger des dizaines d'endroits.

Ce cours commence donc par le geste fondateur du génie logiciel appliqué aux
structures de données : \textbf{séparer ce qu'une structure promet de la
manière dont elle tient sa promesse}. Cette séparation porte un nom --- le
type abstrait de données --- et Java lui donne un support direct : l'interface.
Tout le reste du semestre repose sur ce geste. Quand nous étudierons les
tables de hachage, les arbres ou les graphes, la question sera toujours la
même : quel contrat, et quelle implémentation pour le tenir.

% =====================================================================
\section{Le contrat : type abstrait de données}
% =====================================================================

\subsection{Une promesse, pas une construction}

Quand vous entrez dans une agence de la Poste à Dakar, vous savez comment la
file fonctionne : le premier arrivé est le premier servi, un nouveau client se
met à la fin, et personne ne double. Vous n'avez besoin de savoir ni si les
clients sont debout dans un couloir, ni assis sur des chaises numérotées, ni
inscrits sur un cahier. La règle vous suffit pour vous en servir.

Un type abstrait de données, c'est exactement cela : la règle, sans la mise en
œuvre.

\begin{definition}[Type abstrait de données]
Un \emph{type abstrait de données} (ADT) est la donnée d'un ensemble de
valeurs et d'un ensemble d'opérations sur ces valeurs, spécifiés par leur
\emph{comportement} --- ce qu'elles produisent --- et non par la manière dont
elles sont programmées.
\end{definition}

\begin{definition}[Structure de données]
Une \emph{structure de données} est une réalisation concrète d'un ADT : un
agencement particulier en mémoire, accompagné des algorithmes qui réalisent
chaque opération du contrat.
\end{definition}

Un même ADT admet plusieurs structures de données. Le contrat ne change pas ;
le coût des opérations, lui, change beaucoup. C'est toute la matière de ce
cours.

\begin{figure}[h!]
\centering
\includegraphics[width=0.95\textwidth]{fig_01.pdf}
\caption{Un contrat, trois réalisations. Le code client ne mentionne que
l'interface \texttt{Deque<T>} ; il continue de fonctionner quelle que soit
l'implémentation qu'on lui fournit.}
\end{figure}

\subsection{En Java : l'interface générique}

Java exprime le contrat par une \texttt{interface}, et le rend indépendant du
type des éléments par les \emph{génériques}.

\begin{lstlisting}[caption={Le contrat d'une liste, exprimé en Java}]
public interface Liste<T> {
    void ajouter(T element);        // ajoute en fin
    void ajouter(int i, T element); // insère à la position i
    T get(int i);                   // lit l'élément de rang i
    T supprimer(int i);             // retire et renvoie l'élément de rang i
    int taille();
    boolean estVide();
}
\end{lstlisting}

Le paramètre \texttt{<T>} est un type à fournir à l'usage :
\texttt{Liste<String>} pour des noms, \texttt{Liste<Etudiant>} pour des
objets métier. L'intérêt n'est pas l'élégance : c'est que le compilateur
vérifie le contrat à votre place.

\begin{lstlisting}[caption={Le compilateur garant du contrat}]
Liste<String> noms = new ListeTableau<>();
noms.ajouter("Fatou");
String x = noms.get(0);   // aucun cast : le type est connu
noms.ajouter(42);         // ERREUR À LA COMPILATION, pas à l'exécution
\end{lstlisting}

\begin{appliboxQ}[Mini-exercice 1 --- Contrat ou implémentation ?]
Classez chacune de ces affirmations : relève-t-elle du contrat (ADT) ou de
l'implémentation ?
\begin{enumerate}[nosep,label=(\alph*)]
    \item \texttt{pop()} renvoie le dernier élément empilé.
    \item Les éléments sont stockés dans un tableau de capacité 16.
    \item \texttt{get(i)} lève une exception si $i \geq$ \texttt{taille()}.
    \item Quand le tableau est plein, sa capacité est doublée.
\end{enumerate}
\end{appliboxQ}

\begin{appliboxR}[Réponse]
(a) contrat --- (b) implémentation --- (c) contrat --- (d) implémentation.

\textbf{Le critère :} une clause du contrat est observable par le client à
travers les opérations publiques ; une clause d'implémentation ne l'est pas.
Le doublement de capacité n'est pas observable --- sauf par le temps
d'exécution, et c'est précisément pour cela qu'on devra le mesurer en S2.
\end{appliboxR}

\begin{attentionbox}[Une confusion fréquente]
En Java, \texttt{ArrayList} et \texttt{LinkedList} ne sont pas deux ADTs :
ce sont deux \emph{implémentations} du même ADT, décrit par l'interface
\texttt{java.util.List}. Dire « j'utilise une ArrayList » annonce un choix
technique ; dire « j'ai besoin d'une liste » annonce un besoin. On déclare
toujours le besoin, on n'instancie qu'au dernier moment le choix technique.
\end{attentionbox}

% =====================================================================
\section{List : deux façons de tenir la même promesse}
% =====================================================================

\subsection{Tableau contigu}

Les éléments occupent des cases consécutives. L'adresse de l'élément de rang
$i$ se calcule : \texttt{base} $+\ i \times$ taille d'une case. D'où un accès
direct, en un nombre d'opérations qui ne dépend pas de $i$.

En contrepartie, insérer en position $0$ suppose de décaler tout ce qui suit,
et le tableau a une capacité finie : quand il est plein, il faut en allouer un
plus grand et recopier.

\begin{lstlisting}[caption={Insertion en tête dans un tableau : le décalage}]
public void ajouter(int i, T element) {
    if (taille == donnees.length) agrandir();     // recopie complète
    for (int k = taille; k > i; k--) {
        donnees[k] = donnees[k - 1];              // décalage vers la droite
    }
    donnees[i] = element;
    taille++;
}
\end{lstlisting}

\subsection{Nœuds chaînés}

Chaque élément est logé dans un nœud qui contient, en plus de la valeur, une
référence vers le nœud suivant. Les nœuds sont dispersés en mémoire.
Insérer en tête revient à créer un nœud et à rebrancher une référence : le
travail ne dépend pas du nombre d'éléments déjà présents. Mais atteindre
l'élément de rang $i$ oblige à suivre $i$ références depuis la tête.

\begin{figure}[h!]
\centering
\includegraphics[width=\textwidth]{fig_02.pdf}
\caption{Les deux réalisations de l'ADT List. Ce que l'une donne
gratuitement, l'autre le fait payer.}
\end{figure}

\begin{table}[h!]
\centering
\caption{Coût des opérations de List selon la réalisation}
\begin{tabular}{lcc}
\toprule
\textbf{Opération} & \textbf{Tableau contigu} & \textbf{Nœuds chaînés} \\
\midrule
\texttt{get(i)} & constant & proportionnel à $i$ \\
\texttt{ajouter(x)} en fin & constant (amorti) & constant \\
\texttt{ajouter(0, x)} en tête & proportionnel à $n$ & constant \\
\texttt{supprimer(0)} & proportionnel à $n$ & constant \\
Mémoire par élément & la valeur seule & valeur + référence(s) \\
\bottomrule
\end{tabular}
\end{table}

\begin{appliboxQ}[Mini-exercice 2 --- Choisir sans se tromper]
Un service de la DDD enregistre les passages de bus d'une ligne. Deux usages
distincts :
\begin{enumerate}[nosep,label=(\alph*)]
    \item On ajoute chaque passage à la fin, puis on consulte souvent le
          $k$-ième passage de la journée.
    \item Chaque nouveau passage doit apparaître en tête de liste (le plus
          récent d'abord) et on ne consulte que les premiers.
\end{enumerate}
Quelle réalisation pour chacun ?
\end{appliboxQ}

\begin{appliboxR}[Réponse]
(a) tableau contigu : ajout en fin peu coûteux et accès direct par rang.\\
(b) nœuds chaînés : insertion en tête sans décalage.

\textbf{La leçon :} ce n'est pas la structure qui est « bonne » ou
« mauvaise », c'est l'accord entre la structure et le profil d'usage. On ne
peut donc pas répondre sans connaître les opérations fréquentes.
\end{appliboxR}

% =====================================================================
\section{Stack : le dernier arrivé sort le premier}
% =====================================================================

\subsection{Le contrat}

\begin{lstlisting}[caption={L'ADT Stack}]
public interface Pile<T> {
    void push(T element);   // empile
    T pop();                // dépile et renvoie le sommet
    T peek();               // consulte le sommet sans le retirer
    boolean estVide();
    int taille();
}
\end{lstlisting}

Une seule règle gouverne le tout : \emph{le dernier élément empilé est le
premier dépilé}. Aucune opération ne permet d'atteindre le milieu. Cette
pauvreté délibérée est une force : elle rend l'implémentation efficace et le
raisonnement simple.

\begin{lstlisting}[caption={Implémentation par tableau --- les deux opérations essentielles}]
public class PileTableau<T> implements Pile<T> {
    private Object[] donnees = new Object[16];
    private int taille = 0;

    public void push(T element) {
        if (taille == donnees.length) agrandir();
        donnees[taille] = element;   // on écrit au sommet
        taille++;
    }

    @SuppressWarnings("unchecked")
    public T pop() {
        if (estVide()) throw new IllegalStateException("pile vide");
        taille--;
        T sommet = (T) donnees[taille];
        donnees[taille] = null;      // on libère la référence
        return sommet;
    }
}
\end{lstlisting}

\begin{infobox}[Pourquoi le sommet est en fin de tableau]
Placer le sommet à l'indice \texttt{taille - 1} rend \texttt{push} et
\texttt{pop} sans décalage. Placer le sommet en case $0$ donnerait le même
contrat, mais chaque opération déplacerait tous les éléments. Même contrat,
même résultat, coût incomparable : c'est le thème de la séance prochaine.
\end{infobox}

\subsection{Application : vérifier un parenthésage}

Un éditeur de code doit signaler \texttt{( a + [ b )} comme incorrect. Le
raisonnement est celui d'une pile : chaque symbole ouvrant est mis en attente,
chaque symbole fermant doit correspondre au dernier ouvrant en attente.

\begin{figure}[h!]
\centering
\includegraphics[width=\textwidth]{fig_03.pdf}
\caption{Trace de la pile sur l'expression \texttt{( a + [ b ] )}. La pile
revient vide : l'expression est équilibrée.}
\end{figure}

\begin{lstlisting}[caption={Vérification de parenthésage}]
public static boolean estEquilibree(String expression) {
    Pile<Character> pile = new PileTableau<>();
    for (char c : expression.toCharArray()) {
        if (c == '(' || c == '[' || c == '{') {
            pile.push(c);
        } else if (c == ')' || c == ']' || c == '}') {
            if (pile.estVide()) return false;       // fermante orpheline
            char ouvrante = pile.pop();
            if (!correspond(ouvrante, c)) return false;
        }
    }
    return pile.estVide();   // une ouvrante restée en attente = déséquilibre
}
\end{lstlisting}

\begin{appliboxQ}[Mini-exercice 3 --- Deux façons d'échouer]
Donnez une expression rejetée par le test \texttt{pile.estVide()} de la
ligne~9, et une autre rejetée par le \texttt{return pile.estVide()} final.
\end{appliboxQ}

\begin{appliboxR}[Réponse]
Ligne~9 : \texttt{a + b )} --- une fermante arrive alors que rien n'est en
attente.\\
Retour final : \texttt{( a + b} --- l'analyse se termine avec une ouvrante
encore en attente.

\textbf{Le piège classique :} oublier le test final et ne vérifier que les
correspondances. Le programme accepte alors toute expression à laquelle il
manque des fermantes.
\end{appliboxR}

\subsection{Application : évaluer une expression postfixe}

En notation postfixe, l'opérateur suit ses opérandes : \texttt{3 4 +} vaut
$7$. Aucune parenthèse n'est nécessaire, et l'évaluation se fait en un seul
parcours, à l'aide d'une pile d'opérandes.

\begin{figure}[h!]
\centering
\includegraphics[width=\textwidth]{fig_05.pdf}
\caption{Évaluation de \texttt{3 4 + 2 * 7 -}. Un nombre est empilé ; un
opérateur dépile deux valeurs et empile le résultat. Le résultat final est
le seul élément restant.}
\end{figure}

\begin{lstlisting}[caption={Évaluateur postfixe}]
public static int evaluer(String expression) {
    Pile<Integer> pile = new PileTableau<>();
    for (String jeton : expression.trim().split("\\s+")) {
        switch (jeton) {
            case "+": case "-": case "*": case "/":
                int b = pile.pop();          // ATTENTION à l'ordre
                int a = pile.pop();
                pile.push(appliquer(jeton, a, b));
                break;
            default:
                pile.push(Integer.parseInt(jeton));
        }
    }
    int resultat = pile.pop();
    if (!pile.estVide()) throw new IllegalArgumentException("expression mal formée");
    return resultat;
}
\end{lstlisting}

\begin{appliboxQ}[Mini-exercice 4 --- L'ordre des opérandes]
Que renvoie l'évaluateur sur \texttt{10 3 -} si l'on inverse les deux
dépilements (\texttt{int a = pile.pop(); int b = pile.pop();}) ? Et sur
\texttt{10 3 +} ?
\end{appliboxQ}

\begin{appliboxR}[Réponse]
\texttt{10 3 -} donnerait $3 - 10 = -7$ au lieu de $7$. \texttt{10 3 +}
donnerait $13$, soit la bonne réponse.

\textbf{La leçon :} le bug reste invisible tant qu'on ne teste qu'avec des
opérations commutatives. Vos tests JUnit doivent contenir au moins une
soustraction et une division --- c'est le genre de test qu'on écrit exprès
pour faire échouer une implémentation fausse.
\end{appliboxR}

\begin{appliboxQ}[Mini-exercice 5 --- Trace à la main]
Déroulez l'évaluation de \texttt{5 1 2 + 4 * + 3 -} en notant l'état de la
pile après chaque jeton.
\end{appliboxQ}

\begin{appliboxR}[Réponse]
$[5]$, $[5,1]$, $[5,1,2]$, $[5,3]$, $[5,3,4]$, $[5,12]$, $[17]$, $[17,3]$,
$[14]$. Résultat : $14$.

\textbf{Remarque :} la hauteur maximale atteinte (3) mesure le nombre
d'opérandes simultanément en attente. C'est une quantité utile : elle borne
la mémoire nécessaire.
\end{appliboxR}

% =====================================================================
\section{Queue : le premier arrivé sort le premier}
% =====================================================================

\subsection{Le contrat}

\begin{lstlisting}[caption={L'ADT Queue}]
public interface File<T> {
    void enqueue(T element);  // ajoute en fin
    T dequeue();              // retire et renvoie la tête
    T peek();                 // consulte la tête
    boolean estVide();
    int taille();
}
\end{lstlisting}

\subsection{Le tableau circulaire}

L'implémentation naïve --- tête en case $0$, décalage à chaque
\texttt{dequeue} --- tient le contrat mais paie un décalage complet à chaque
retrait. On évite ce coût en ne déplaçant plus les éléments, mais l'indice de
la tête.

Deux variables suffisent : \texttt{tete} (indice du prochain élément à
servir) et \texttt{taille}. La fin de la file se calcule :
\[
\texttt{fin} \;=\; (\texttt{tete} + \texttt{taille}) \bmod m
\]
où $m$ est la capacité du tableau. Le modulo fait « repasser » la fin au
début du tableau quand elle atteint le bord : le tableau se comporte comme un
anneau.

\begin{figure}[h!]
\centering
\includegraphics[width=\textwidth]{fig_04.pdf}
\caption{Un guichet, six cases. À droite, la file occupe les cases $2$ à $5$
puis reprend en case $0$ : les éléments n'ont jamais été déplacés.}
\end{figure}

\begin{lstlisting}[caption={File circulaire : les deux opérations}]
public class FileCirculaire<T> implements File<T> {
    private Object[] donnees = new Object[8];
    private int tete = 0;
    private int taille = 0;

    public void enqueue(T element) {
        if (taille == donnees.length) agrandir();
        int fin = (tete + taille) % donnees.length;
        donnees[fin] = element;
        taille++;
    }

    @SuppressWarnings("unchecked")
    public T dequeue() {
        if (estVide()) throw new IllegalStateException("file vide");
        T premier = (T) donnees[tete];
        donnees[tete] = null;
        tete = (tete + 1) % donnees.length;   // la tête avance, rien ne bouge
        taille--;
        return premier;
    }
}
\end{lstlisting}

\begin{appliboxQ}[Mini-exercice 6 --- Arithmétique de l'anneau]
Capacité $m = 6$, \texttt{tete} $= 4$, \texttt{taille} $= 3$.
\begin{enumerate}[nosep,label=(\alph*)]
    \item Quelles cases sont occupées ?
    \item Dans quelle case écrit le prochain \texttt{enqueue} ?
    \item Après deux \texttt{dequeue}, que vaut \texttt{tete} ?
\end{enumerate}
\end{appliboxQ}

\begin{appliboxR}[Réponse]
(a) cases $4$, $5$, $0$ --- car $(4+0) \bmod 6 = 4$, $(4+1) \bmod 6 = 5$,
$(4+2) \bmod 6 = 0$.\\
(b) $(4 + 3) \bmod 6 = 1$.\\
(c) $\texttt{tete} = (4 + 2) \bmod 6 = 0$, et \texttt{taille} $= 1$.

\textbf{Le piège classique :} utiliser \texttt{tete < fin} pour tester si la
file est vide. C'est faux dès que la file a « bouclé ». La seule information
fiable est \texttt{taille}.
\end{appliboxR}

\begin{attentionbox}[Pourquoi on garde \texttt{taille} et pas seulement deux indices]
Avec les seuls indices \texttt{tete} et \texttt{fin}, la file vide et la file
pleine donnent la même configuration \texttt{tete == fin} : impossible de les
distinguer. Deux remèdes existent : maintenir un compteur \texttt{taille}
(retenu ici, le plus lisible), ou sacrifier une case du tableau.
\end{attentionbox}

% =====================================================================
\section{Deque : les deux extrémités}
% =====================================================================

Un \emph{deque} (\emph{double-ended queue}) accepte ajouts et retraits aux
deux bouts. Il subsume la pile et la file : un deque restreint à
\texttt{addLast}/\texttt{removeLast} se comporte comme une pile, restreint à
\texttt{addLast}/\texttt{removeFirst} comme une file.

\begin{lstlisting}[caption={L'ADT Deque --- le contrat que vous implémenterez en Lab}]
public interface Deque<T> {
    void addFirst(T element);
    void addLast(T element);
    T removeFirst();
    T removeLast();
    T peekFirst();
    T peekLast();
    int size();
    boolean isEmpty();
}
\end{lstlisting}

L'implémentation par tableau circulaire redimensionnable rend les huit
opérations peu coûteuses, car aucune ne déplace les éléments : \texttt{addFirst}
recule l'indice de tête modulo $m$, \texttt{addLast} avance l'indice de fin.

\begin{lstlisting}[caption={Reculer un indice sans nombre négatif}]
// FAUX si tete vaut 0 : (0 - 1) % m donne -1 en Java
int nouvelleTete = (tete - 1) % donnees.length;

// CORRECT : on ajoute m avant de prendre le modulo
int nouvelleTete = (tete - 1 + donnees.length) % donnees.length;
\end{lstlisting}

\begin{appliboxQ}[Mini-exercice 7 --- Le modulo négatif]
En Java, que vaut \texttt{(0 - 1) \% 8} ? Et \texttt{(0 - 1 + 8) \% 8} ?
\end{appliboxQ}

\begin{appliboxR}[Réponse]
$-1$ et $7$.

\textbf{Pourquoi :} en Java, l'opérateur \texttt{\%} est un reste, pas un
modulo mathématique : son signe est celui du dividende. Un indice $-1$
provoque une \texttt{ArrayIndexOutOfBoundsException} à la première écriture.
C'est l'erreur la plus fréquente de l'implémentation de \texttt{ArrayDeque} :
notez-la, vous la rencontrerez en Lab.
\end{appliboxR}

% =====================================================================
\section{Choisir}
% =====================================================================

La démarche d'ingénieur s'énonce en trois temps, dans cet ordre.

\begin{enumerate}
    \item \textbf{Quelles opérations le problème demande-t-il ?} On identifie
          l'ADT. Un annuleur d'actions (\texttt{Ctrl+Z}) demande le dernier
          geste effectué : c'est une pile. Un guichet sert dans l'ordre
          d'arrivée : c'est une file.
    \item \textbf{Quelles opérations sont fréquentes ?} On choisit la
          structure concrète selon le profil d'usage.
    \item \textbf{Qu'observe-t-on en exécution ?} On mesure, et l'on révise
          si la mesure contredit la prédiction. C'est l'objet de S2.
\end{enumerate}

\begin{table}[h!]
\centering
\caption{Repères pour le choix de l'ADT}
\begin{tabular}{lll}
\toprule
\textbf{Besoin exprimé} & \textbf{ADT} & \textbf{Réalisation courante} \\
\midrule
Accès par rang, parcours ordonné & List & tableau contigu \\
Insertions/retraits en tête fréquents & List & nœuds chaînés \\
Annulation, retour en arrière & Stack & tableau (sommet en fin) \\
Service dans l'ordre d'arrivée & Queue & tableau circulaire \\
Ajout/retrait aux deux bouts & Deque & tableau circulaire \\
\bottomrule
\end{tabular}
\end{table}

\begin{appliboxQ}[Mini-exercice 8 --- Du besoin à l'ADT]
Nommez l'ADT adapté à chacun de ces besoins :
\begin{enumerate}[nosep,label=(\alph*)]
    \item Le navigateur revient à la page précédente.
    \item Une imprimante partagée traite les travaux dans l'ordre reçu.
    \item Un historique où l'on consulte le plus récent, mais où l'on purge
          le plus ancien quand la capacité est atteinte.
\end{enumerate}
\end{appliboxQ}

\begin{appliboxR}[Réponse]
(a) Stack --- (b) Queue --- (c) Deque : on lit et écrit d'un côté (le
récent), on purge de l'autre (l'ancien).

\textbf{Remarque :} (c) est exactement le mécanisme d'un cache à éviction du
plus ancien. On retrouvera ce motif quand on parlera de tables de hachage et
de files de priorité.
\end{appliboxR}

\begin{appliboxQ}[Mini-exercice 9 --- Le coût caché]
Un étudiant écrit une file à l'aide d'une \texttt{ArrayList} :
\texttt{enqueue} fait \texttt{liste.add(x)} et \texttt{dequeue} fait
\texttt{liste.remove(0)}. Le contrat est-il respecté ? Le choix est-il bon ?
\end{appliboxQ}

\begin{appliboxR}[Réponse]
Le contrat est respecté : les éléments sortent bien dans l'ordre d'arrivée.
Le choix, lui, est mauvais : \texttt{remove(0)} décale tous les éléments
restants, donc chaque retrait coûte proportionnellement au nombre d'éléments
présents.

\textbf{La leçon centrale de la séance :} un programme peut être
\emph{correct} et \emph{inefficace}. La correction se vérifie par des tests ;
l'efficacité demande un autre outil --- c'est pourquoi S2 existe.
\end{appliboxR}

% =====================================================================
\begin{retenirbox}[Points clés à retenir]
\begin{enumerate}[nosep]
    \item Un ADT est un contrat : un ensemble d'opérations spécifiées par
          leur comportement. Une structure de données est une manière de
          tenir ce contrat.
    \item En Java, le contrat s'écrit \texttt{interface}, et les génériques
          le rendent indépendant du type des éléments. On déclare avec le
          type de l'interface, on n'instancie qu'au dernier moment.
    \item Tableau contigu et nœuds chaînés tiennent le même contrat List avec
          des coûts opposés : accès direct contre insertion en tête.
    \item La pile obéit au dernier-arrivé-premier-sorti ; elle résout la
          vérification de parenthésage et l'évaluation postfixe en un seul
          parcours.
    \item La file obéit au premier-arrivé-premier-sorti ; le tableau
          circulaire lui évite tout décalage, au prix d'une arithmétique
          modulo.
    \item En Java, \texttt{\%} est un reste et peut être négatif : pour
          reculer un indice, écrire \texttt{(i - 1 + m) \% m}.
    \item Le deque généralise pile et file ; c'est la structure que vous
          implémenterez au Lab et dans le Projet P1.
\end{enumerate}
\end{retenirbox}

% =====================================================================
\section*{Ce qui vient ensuite : Analyse algorithmique}
\addcontentsline{toc}{section}{Ce qui vient ensuite : Analyse algorithmique}
% =====================================================================

\begin{prereqbox}[Vers la séance suivante]
\textbf{Ce qu'on sait faire maintenant :} énoncer un contrat sous forme
d'interface générique, et le réaliser de plusieurs manières --- tableau
contigu, nœuds chaînés, tableau circulaire.\\[4pt]
\textbf{La limite qu'on va dépasser :} tout au long de cette séance, nous
avons dit « coûteux », « constant », « proportionnel à $n$ » sans jamais rien
démontrer. Sur quoi repose l'affirmation que \texttt{remove(0)} est mauvais ?
Comment comparer deux implémentations autrement qu'à l'intuition ?\\[4pt]
\textbf{En S2 :} on compte les opérations élémentaires, on définit
formellement $\bigO$, $\bigOmega$ et $\bigTheta$, on prouve des bornes, et on
apprend à analyser le code récursif. Les affirmations de cette séance
deviendront des énoncés démontrés.
\end{prereqbox}

% =====================================================================
\section*{Références}
\addcontentsline{toc}{section}{Références}
% =====================================================================

\begin{itemize}[nosep]
    \item R. Sedgewick, K. Wayne, \emph{Algorithms}, 4\textsuperscript{e} éd.
          --- §1.2 (Data Abstraction), §1.3 (Bags, Queues, and Stacks).
    \item T. Cormen, C. Leiserson, R. Rivest, C. Stein,
          \emph{Introduction to Algorithms}, 4\textsuperscript{e} éd.
          --- Ch.~10 (Elementary Data Structures).
    \item S. Skiena, \emph{The Algorithm Design Manual},
          3\textsuperscript{e} éd. --- Ch.~3 (Data Structures).
\end{itemize}

\end{document}
